\begin{theorem}[一阶局部作用量的 Euler--Lagrange 方程]
对 $S[\phi]=\int\mathcal L(\phi,\partial_\mu\phi,x)\,\dd^dx$，若变分在边界消失，则驻值条件等价于 $\partial\mathcal L/\partial\phi_a-\partial_\mu(\partial\mathcal L/\partial(\partial_\mu\phi_a))=0$。复场的共轭分量按独立实自由度处理。
\end{theorem}

\begin{proof}
正文\eqnrefs{eq:field-action-general-dp61} 把场的整段时空历史映射成作用量。对场做任意小变分时，唯一需要处理的是含 $\partial_\mu\delta\Phi_A$ 的项；把它分部积分后，边界条件 $\delta\Phi_A|_{\partial\Omega}=0$ 消去表面项，而区域内部的 $\delta\Phi_A(x)$ 又可以独立任取，因此体积分系数必须逐点为零。这正给出正文\eqnrefs{eq:field-euler-lagrange-dp61}。

从定义出发，
\begin{align*}
 \delta S
 &=\frac1c\int_\Omega\dd^4 x\,
 \left[
 \frac{\partial\Lag}{\partial\Phi_A}\delta\Phi_A
 +\frac{\partial\Lag}{\partial(\partial_\mu\Phi_A)}
 \delta(\partial_\mu\Phi_A)
 \right],\\
 &=\frac1c\int_\Omega\dd^4 x\,
 \left[
 \frac{\partial\Lag}{\partial\Phi_A}\delta\Phi_A
 +\frac{\partial\Lag}{\partial(\partial_\mu\Phi_A)}
 \partial_\mu\delta\Phi_A
 \right].
\end{align*}
第二项使用乘积求导反解：
\begin{align*}
 \frac{\partial\Lag}{\partial(\partial_\mu\Phi_A)}
 \partial_\mu\delta\Phi_A
 &=\partial_\mu\!\left[
 \frac{\partial\Lag}{\partial(\partial_\mu\Phi_A)}
 \delta\Phi_A\right]
 -\partial_\mu\!\left[
 \frac{\partial\Lag}{\partial(\partial_\mu\Phi_A)}\right]
 \delta\Phi_A.
\end{align*}
代回以后，
\begin{align*}
 \delta S
 &=\frac1c\int_\Omega\dd^4 x\,
 \left[
 \frac{\partial\Lag}{\partial\Phi_A}
 -\partial_\mu\frac{\partial\Lag}{\partial(\partial_\mu\Phi_A)}
 \right]\delta\Phi_A\\
 &\quad+\frac1c\int_\Omega\dd^4 x\,
 \partial_\mu\!\left[
 \frac{\partial\Lag}{\partial(\partial_\mu\Phi_A)}
 \delta\Phi_A\right].
\end{align*}
四维 Gauss 定理把最后一项变成 $\partial\Omega$ 上的面积分；由于边界上 $\delta\Phi_A=0$，该项消失。于是
\begin{equation*}
 \delta S
 =\frac1c\int_\Omega\dd^4 x\,
 \mathfrak E_A(x)\,\delta\Phi_A(x),
\end{equation*}
其中
\begin{equation*}
 \mathfrak E_A
 \equiv
 \frac{\partial\Lag}{\partial\Phi_A}
 -\partial_\mu\frac{\partial\Lag}{\partial(\partial_\mu\Phi_A)}.
\end{equation*}
区域内部的 $\delta\Phi_A(x)$ 可以独立任取；要使 $\delta S=0$ 对所有变分都成立，只能有 $\mathfrak E_A(x)=0$。移项即得到\eqnrefs{eq:field-euler-lagrange-dp61}。
\end{proof}

\begin{theorem}[Noether 流的在壳守恒]
若无穷小场变换使拉氏量密度变化为全导数 $\delta\mathcal L=\partial_\mu K^\mu$，则 $j^\mu=\frac{\partial\mathcal L}{\partial(\partial_\mu\phi_a)}\delta\phi_a-K^\mu$ 在 Euler--Lagrange 方程成立时满足 $\partial_\mu j^\mu=0$；坐标变换需使用本章的固定坐标总变分。
\end{theorem}

\begin{proof}
链式法则给

\begin{align*}
 \delta\Lag
 &=\frac{\partial\Lag}{\partial\Phi_A}\delta\Phi_A
 +\frac{\partial\Lag}{\partial(\partial_\mu\Phi_A)}
 \partial_\mu\delta\Phi_A.
\end{align*}
对第二项做同样的乘积拆分，
\begin{align*}
 \delta\Lag
 &=\left[
 \frac{\partial\Lag}{\partial\Phi_A}
 -\partial_\mu\frac{\partial\Lag}{\partial(\partial_\mu\Phi_A)}
 \right]\delta\Phi_A\\
 &\quad+\partial_\mu\!\left[
 \frac{\partial\Lag}{\partial(\partial_\mu\Phi_A)}
 \delta\Phi_A\right].
\end{align*}
在满足 Euler--Lagrange 方程的经典解上，第一行的方括号为零。再代入 $\delta\Phi_A=\epsilon\Delta_A$：
\begin{equation*}
 \delta\Lag
 =\epsilon\,\partial_\mu\!\left[
 \sum_A\frac{\partial\Lag}{\partial(\partial_\mu\Phi_A)}
 \Delta_A\right].
\end{equation*}
另一方面，对称性假设给 $\delta\Lag=\epsilon\partial_\mu K^\mu$。两式相减，并利用常数 $\epsilon$ 可约去，得到
\begin{equation*}
 \partial_\mu\!\left[
 \sum_A\frac{\partial\Lag}{\partial(\partial_\mu\Phi_A)}
 \Delta_A-K^\mu\right]=0,
\end{equation*}
即\eqnrefs{eq:noether-current-general-dp61}。
\end{proof}
